\chapter{Despliegue}
\label{cap:despliegue}

\section{Requisitos del sistema}

Para desplegar la plataforma Atlas se requiere una base mínima de infraestructura y versiones de \textit{runtime} coherentes entre entornos.

\begin{itemize}
\item Node.js 20 o superior para la ejecución de microservicios y utilidades de build.
\item Docker 24 o superior para empaquetado y despliegue de contenedores.
\item Cluster Kubernetes (on-premise y/o cloud) para orquestación de pods stateless.
\item RabbitMQ 3.x como bus de mensajería AMQP (incluyendo instancia IoT con MQTT en CPD).
\item MariaDB 11.x para persistencia relacional.
\item Valkey 8.x para caché y almacenamiento temporal.
\item Ollama (on-premise) para capacidades LLM, OCR y embeddings en el puerto \texttt{11434}.
\item NGINX para publicación HTTP/HTTPS y terminación TLS.
\end{itemize}

En entornos de producción se recomienda separar credenciales, certificados y configuración por dominio funcional (API pública, CMS, servicios internos e IoT).

\section{Entorno de desarrollo}

El entorno local está orientado a validar cambios funcionales y de integración antes de promocionarlos a preproducción o producción.

Flujo recomendado:

\begin{enumerate}
\item Clonar los repositorios del ecosistema Atlas (\texttt{atlas-*} y \texttt{ferimer-*}) en el directorio de trabajo.
\item Configurar variables de entorno para cada servicio (RabbitMQ, MariaDB, Valkey y endpoints internos).
\item Levantar servicios de infraestructura compartida (RabbitMQ, MariaDB, Valkey y, cuando aplique, Ollama).
\item Iniciar los microservicios de aplicación (API, CMS, scraping y servicios IA).
\item Verificar conectividad extremo a extremo ejecutando una llamada API que dispare una operación RPC sobre RabbitMQ.
\item Revisar trazas y \textit{correlation id} para confirmar que el flujo de petición/respuesta es correcto.
\end{enumerate}

La validación local debe incluir escenarios síncronos (RPC) y asíncronos (Fire-and-Forget), así como pruebas básicas de esquemas JSON en petición y respuesta.

\section{Entorno de producción}

En producción se despliegan contenedores inmutables sobre Kubernetes, manteniendo separación entre entorno on-premise y cloud.

Pautas operativas:

\begin{itemize}
\item Construcción de imágenes versionadas por servicio y promoción por entornos.
\item Inyección de configuración mediante variables de entorno y \textit{ConfigMaps}.
\item Gestión de secretos (credenciales, tokens, certificados) en mecanismos de \textit{secrets} del cluster.
\item Definición de \textit{liveness/readiness probes} para facilitar reinicio automático y balanceo saludable.
\item Estrategia de despliegue progresivo (\textit{rolling update}) para minimizar indisponibilidad.
\item Observabilidad centralizada con métricas, logs y trazabilidad por \textit{correlation id}.
\end{itemize}

Cuando se utilicen composiciones tipo Docker Compose para entornos de integración o staging, deben tratarse como apoyo de validación, no como mecanismo principal de operación productiva.

\section{Configuración de los servicios}

La configuración se organiza por servicio y por entorno. Como principio general, los valores por defecto residen en ficheros de configuración versionados (por ejemplo, \texttt{config/default.cjs}) y los valores sensibles se externalizan por variables de entorno.

Variables habituales:

\begin{itemize}
\item Conectividad AMQP: host, puerto, vhost (\texttt{/atlas}), usuario, contraseña.
\item Persistencia: host/puerto de MariaDB, nombre de base de datos y credenciales.
\item Caché: endpoint de Valkey y parámetros de expiración.
\item IA: endpoint de Ollama, modelos por defecto y límites de timeout.
\item Publicación: puertos HTTP/HTTPS y dominios para NGINX/API/CMS.
\end{itemize}

Se recomienda mantener un inventario de variables por repositorio, junto con ejemplos de \texttt{.env} sin secretos, para facilitar despliegues reproducibles.

\section{Publicación de librerías}

Las librerías compartidas del ecosistema (especialmente \texttt{ferimer-*}) se publican en GitHub Packages con versionado semántico.

Proceso recomendado:

\begin{enumerate}
\item Incrementar versión siguiendo SemVer (\textit{patch}, \textit{minor} o \textit{major}).
\item Generar \textit{changelog} técnico con cambios funcionales y de compatibilidad.
\item Publicar paquete en el registro privado de GitHub Packages.
\item Consumir la nueva versión desde servicios dependientes y ejecutar batería de regresión.
\end{enumerate}

La autenticación para publicación/consumo se realiza mediante \textit{Personal Access Token} (PAT) con permisos mínimos de lectura/escritura de paquetes según rol.
